% =============================================================================
% 017 / 068
% Status: raw
% =============================================================================
% Notes:
%
% Source question:
% 母語詞彙可以從別的語言借用，這叫做外來語，
% 這是很常見的現象。
% 
% 那構詞法和文法也可以從別的語言借過來，並紮根於母語中嗎？
% =============================================================================

\chapter[母語詞彙可以從別的語言借用，這叫做外來語， 這是很常見的現象。 那構詞法和文法也可以從別的語言借過來，並紮根於母語中嗎？]{母語詞彙可以從別的語言借用，這叫做外來語， \\[0.35em] 這是很常見的現象。 \\[0.35em] 那構詞法和文法也可以從別的語言借過來，並紮根於母語中嗎？}
\qaanswer{
最穩妥的方式是按照經驗主義的方式來構建，也就是說外來語的文法、構詞法首先要能被吳語區的人在使用時普遍接受，但同時又不能與吳語本有的文法、構詞法產生邏輯矛盾。比如吳語非常講究區別詞性，而漢語就很難區分詞性了。基本上是中古漢語、近古漢語、梵語、中古波斯語、契丹語的一些詞彙作為詞根融入到了吳語中。直接從外來語拿來構詞法和文法我相信很難有人認可，也不知如何來使用他。現在主要的問題是吳語區的人掌握的吳語基本也都是些被官話、普通話嚴重滲透失去吳語語法特點的藍青官話。按我的標準這些都算不上吳語，是一種帶著吳語口音、偶爾有吳語詞彙的「普通話」。這些人對於吳語語法的瞭解、構詞法的瞭解基本上是一竅不通的。建構的前提是瞭解自身。與其一開始考慮借鑑別人不如先好好了解自己。懂得自身才能知道外界哪些是需要的哪些是不需要的。
}

